\documentclass[aspectratio=169,10pt]{beamer}
\usetheme[numbering=fraction,progressbar=frametitle]{metropolis}
\usepackage{fontspec}
\setsansfont{Latin Modern Sans}
\usepackage{amsmath,amssymb,mathtools,graphicx,booktabs}
\usepackage{polyglossia}
\setdefaultlanguage{french}
\setbeamertemplate{navigation symbols}{}
\setbeamerfont{frametitle}{size=\large}
\setbeamerfont{normal text}{size=\normalsize}
\setlength{\parskip}{.45em}
\hypersetup{pdfauthor={},pdfsubject={Lecture slides},pdfkeywords={},pdftitle={Boucles et organigrammes}}
\providecommand{\N}{\mathbb{N}}
\providecommand{\Z}{\mathbb{Z}}
\providecommand{\D}{\mathbb{D}}
\providecommand{\Q}{\mathbb{Q}}
\providecommand{\R}{\mathbb{R}}
\providecommand{\C}{\mathbb{C}}
\providecommand{\vect}[1]{\overrightarrow{#1}}
\providecommand{\vi}{\vec{\imath}}
\providecommand{\vj}{\vec{\jmath}}
\providecommand{\vk}{\vec{k}}
\providecommand{\norm}[1]{\left\lVert #1\right\rVert}
\providecommand{\vc}[2]{\begin{pmatrix}#1\\#2\end{pmatrix}}
\providecommand{\deter}[2]{\det\left(#1,#2\right)}
\providecommand{\abs}[1]{\left\lvert #1\right\rvert}
\providecommand{\dd}{\,\mathrm{d}}
\providecommand{\e}{\mathrm{e}}
\providecommand{\Card}{\operatorname{Card}}
\providecommand{\E}{\operatorname{E}}
\providecommand{\V}{\operatorname{V}}
\providecommand{\degre}{^{\circ}}
\providecommand{\un}[1]{\,\mathrm{#1}}
\providecommand{\sci}[2]{#1\times10^{#2}}
\providecommand{\pgcd}{\operatorname{PGCD}}
\providecommand{\ppcm}{\operatorname{PPCM}}
\title{Boucles et organigrammes}
\subtitle{Cours et exemples}
\date{}
\author{}
\begin{document}
\begin{frame}[plain]
\vfill
{\usebeamerfont{title}\inserttitle\par}
\medskip
{\usebeamerfont{subtitle}\insertsubtitle\par}
\medskip
{\usebeamercolor[fg]{progress bar}\rule{\textwidth}{.5pt}}
\vfill
\end{frame}
\begin{frame}{Objectifs}
Exécuter et construire une boucle, produire un organigramme et justifier la terminaison d’un algorithme simple.
\end{frame}
\begin{frame}{Exécuter sans sauter une instruction}
Une affectation \(a\leftarrow E\) remplace la valeur de \(a\) par la valeur actuelle de \(E\). Elle ne constitue pas une égalité à résoudre. Une boucle « pour » parcourt une liste d’indices; une boucle « tant que » répète tant que sa condition est vraie. Le test de « tant que » précède chaque tour, y compris le premier.
\end{frame}
\begin{frame}{Exécuter sans sauter une instruction}
\begin{center}\begin{tabular}{l}
Lire un entier \(n\ge0\) \\ 
\(S\leftarrow0\) \\ 
Pour \(k\) allant de \(1\) à \(n\) \\ 
  \(S\leftarrow S+k\) \\ 
Fin pour \\ 
Afficher \(S\) \\ 
\end{tabular}\end{center}
\end{frame}
\begin{frame}{Exécuter sans sauter une instruction (2)}
Pour \(n=4\), les valeurs successives de \(S\) sont \(0,1,3,6,10\). Si \(n=0\), aucun tour n’est effectué et la somme vide vaut \(0\).
\end{frame}
\begin{frame}{Exécuter sans sauter une instruction}
Après le tour d’indice \(k\), \(S=1+\cdots+k\).
\end{frame}
\begin{frame}{Exécuter sans sauter une instruction}
Avant le premier tour, \(S=0\). Si la somme est correcte au tour précédent, l’ajout de \(k\) donne celle des \(k\) premiers entiers. Une récurrence justifie chaque tour; après \(n\) tours, le résultat annoncé est obtenu.
\end{frame}
\begin{frame}{Boucle conditionnelle et terminaison}
Avec \(u\leftarrow1\), \(n\leftarrow0\), tant que \(u<20\), faire \(u\leftarrow2u\) puis \(n\leftarrow n+1\). Les valeurs de \(u\) sont \(1,2,4,8,16,32\); la sortie est \(n=5\). Après \(n\) tours, \(u=2^n\). La suite des puissances dépasse \(20\) : l’arrêt est justifié.
\end{frame}
\begin{frame}{Boucle conditionnelle et terminaison}
Si le corps ne modifie pas ce qui intervient dans le test, la boucle peut être infinie. Pour permuter deux valeurs, \(a\leftarrow b\) puis \(b\leftarrow a\) ne convient pas : la première valeur de \(a\) a été perdue. Utiliser une variable temporaire.
\end{frame}
\begin{frame}{Boucle conditionnelle et terminaison}
Préciser les entrées autorisées; dresser une table d’exécution; expliquer ce qui reste vrai après chaque tour (invariant); expliquer pourquoi le test finit par être faux (terminaison); interpréter la sortie et les cas limites.
\end{frame}
\begin{frame}{Boucle conditionnelle et terminaison}
On initialise \(a=7,b=3\), puis on exécute \(a\leftarrow a-b\) et \(b\leftarrow a+b\).

\begin{center}\begin{tabular}{lll}
État & \(a\) & \(b\) \\ 
Initial & 7 & 3 \\ 
Après la première affectation & 4 & 3 \\ 
Après la seconde & 4 & 7 \\ 
\end{tabular}\end{center}
\end{frame}
\begin{frame}{Boucle conditionnelle et terminaison (2)}
Dans la seconde ligne de calcul, \(a\) vaut déjà \(4\). Remplacer les deux variables simultanément par leurs anciennes valeurs décrirait un autre algorithme.
\end{frame}
\begin{frame}{Boucle conditionnelle et terminaison}
Un invariant explique ce que contiennent les variables si l’algorithme atteint une étape. Il ne garantit pas que l’algorithme s’arrête. Pour une boucle bornée, compter les tours; pour une boucle conditionnelle, trouver une quantité qui finit par franchir le seuil. Tester aussi le cas où la condition est fausse dès le départ.
\end{frame}
\begin{frame}{Boucle conditionnelle et terminaison}
Pour calculer \(u_0+\cdots+u_N\) quand \(u_{n+1}=2u_n+1\), initialiser \(u=u_0,S=u_0\). Au tour \(k=1,\ldots,N\), mettre d’abord \(u\leftarrow2u+1\), puis \(S\leftarrow S+u\). Après le tour \(k\), \(u=u_k\) et \(S=u_0+\cdots+u_k\). En \(N=0\), il n’y a pas de tour et la sortie est déjà correcte. Inverser les deux mises à jour compterait deux fois \(u_0\) et oublierait le dernier terme.
\end{frame}
\begin{frame}{Boucle conditionnelle et terminaison}
Une suite peut tendre vers un seuil sans l’atteindre à un rang fini. Un test d’égalité exacte ne doit pas être remplacé tacitement par un test approché. Pour demander une précision, fixer une tolérance strictement positive dans l’énoncé.
\end{frame}
\end{document}
